\section{Related work}
\label{sec:related}

\paragraph{Channel pruning.}
Most pipelines start from one of three scores: the $\ell_1$ norm of a filter~\cite{li2017pruning}, the BatchNorm scale learned under a sparsity penalty~\cite{liu2017slimming}, or a first-order Taylor estimate of the loss change~\cite{molchanov2017pruning,molchanov2019importance}. He et al.~\cite{he2017channel} choose channels that best reconstruct the next layer's output and, for residual networks, prune only the inner convolutions of each block, which is the scheme we use. Discrimination-aware pruning~\cite{zhuang2018dcp} adds auxiliary classification losses to intermediate layers and keeps the channels that preserve them. DACIS asks a similar question with a cheaper statistic, the Fisher ratio of the pooled activation, and no extra training. These methods are evaluated on the classes they were trained for. The channels we keep are judged on classes the network never saw. MetaPruning~\cite{liu2019metapruning} also uses the word ``meta'', but there a hypernetwork predicts weights for candidate architectures; our meta stage is ordinary episodic training whose gradients feed a later pruning step.

\paragraph{Few-shot classification.}
Prototypical networks~\cite{snell2017proto} label a query by the nearest class mean in embedding space, and MAML~\cite{finn2017maml} learns an initialisation that adapts in a few gradient steps. Two findings shape our protocol. A network pre-trained with plain cross-entropy on the base classes is a strong baseline~\cite{chen2019closer,tian2020rethinking}, and episodic training on top of it with a cosine nearest-centroid head can add a little more~\cite{chen2021metabaseline}. Our stages S0 and M follow that recipe. The cross-domain benchmark of Guo et al.~\cite{guo2020bscd} has crop-disease and EuroSAT targets and reports that meta-learning methods often lose to simple fine-tuning when the target is far from the source. Section~\ref{sec:results} finds the same sign on both domains.

\paragraph{Plants and Earth observation.}
PlantVillage~\cite{hughes2015plantvillage} photographs single leaves against a plain background. PlantDoc~\cite{singh2020plantdoc} was assembled from internet photographs because large non-laboratory plant datasets were missing, and Barbedo~\cite{barbedo2018factors} lists image background and capture conditions among the factors that limit deep-learning disease recognition. Argüeso et al.~\cite{argueso2020fewshot} study few-shot recognition on PlantVillage with a $32/6$ class split and a Siamese network; we use a similar cross-class split with prototype classification. In remote sensing, few-shot scene classification has been approached with metric learning on aerial benchmarks~\cite{li2021rsmetanet}, and the Sentinel-2~\cite{drusch2012sentinel2} patches of EuroSAT~\cite{helber2019eurosat} are a standard small benchmark. On-board inference is not hypothetical: CloudScout~\cite{giuffrida2020cloudscout} runs a compact CNN on a Myriad~2 accelerator in the $\Phi$-Sat-1 mission~\cite{giuffrida2022phisat}. We do not target that hardware. Our question is narrower: when a ResNet-18 must shrink and must still recognise classes it was not trained on, does the choice of channels matter?

\paragraph{Earlier version of this method.}
A preprint~\cite{uddin2026preprint} introduced the names PMP and DACIS and reported results from an earlier implementation. We rebuilt both from scratch, changed the protocol (class-disjoint splits, shared depth profile, paired episodes, five partitions), and reuse none of its numbers.
